

\section{Conclusion}
\label{sec:rh_conclusion}
\begin{itemize}
\item \textbf{Exogenous predictors in intraday realized variance forecasts.} We ask what exogenous predictors, from
  the futures' own realized moments and the constituents' trading to the FOMC calendar, implied volatility and
  alternative data, add to forecasts of the E-mini S\&P~500's realized variance on every regular-hours bar,
  against OLS on the HAR benchmark set.
\item \textbf{The exogenous predictors improve the forecast of every bar} (Section~\ref{sec:rh_results}).
\item \textbf{The gain is concentrated} (Section~\ref{sec:rh_fi}). The FOMC calendar
  carries most of the trees' gain and the futures' realized moments come second. Along the feature-count ladder the
  \ladderEBTopK{} highest-ranked series retain \ladderEBTreeRetainRng{}\% of the boosted trees' gain and
  \rhldEBLinRTenMin{}--\rhldEBLinRTenMax{}\% of the linear models'. The gain sits in about \ladderEBTopK{} of the
  \rhldEBNSeries{} series rather than spreading thinly over many.
\item \textbf{Beyond the two groups, a small gain in the boosted trees.}
\item \textbf{QLIKE and squared error.} Under squared-error loss the full set is not more accurate (its mean squared
  error is higher in \rhcaNMseHigherBarSpan{} of \rhorAfPairs{} classes, significantly in \rhcaNMseHigherSigBarSpan),
  a difference carried almost entirely by \rhcaTopDays{} days of the COVID-19 crash; without the crash window it is
  significantly lower in \rhcaNMseLowerSigBarOosExCrash{} classes.
\item \textbf{An application to 0DTE straddles} (Section~\ref{sec:rh_econ}).
\item \textbf{Comparison with machine-learning option portfolios.}
\item \textbf{Scope and extensions.} One index and one futures contract; a window fixed after an out-of-sample
  comparison of windows, all reported in the online appendix; NYSE-holiday sessions kept in the estimation sample; a
  short option sample (\rheSzDays{} days). Extensions: the sessions after \rhorSpanEnd, other markets, density
  forecasts, and trades entered at earlier bars from the every-bar forecasts already produced.
\end{itemize}


